\section{Related Work}

% \begin{table*}[]
%     \centering
%     \begin{tabular}{|c|c|c|c|}
%     \hline
%          & inference & decomposition & generality \\
%          \hline\hline
%          Prompt 1 & high & holistic/analytic & general \\
%          Prompt 4 & high/low & holistic/analytic & general \\
%          Prompt 6 & low & holistic/analytic & task specific/general \\
%          \hline
%          IF & low/high or low & holistic/analytic & general/task specific \\
%          \hline
%     \end{tabular}
%     \caption{Components of prompts}
%     \label{tab:levels}
% \end{table*}

\subsection{Rubrics in Context}
% \jessica{talk about rubric terms, maybe take from the intro paragraph of experiment setup - decomposition level, generality level, agreement; how we can view these in the context of ML "rubrics"}

Expanding upon Section~\ref{sec:intro}, rubrics consist of both the criteria, which are the components of the overall evaluation, and the descriptions for the criteria. 

\textbf{Decomposition level} describes criteria presentation, and refers to whether prompts are \textit{holistic}, where ``all criteria are [applied] at the same time'', or \textit{analytic}, where ``work [is described] on each criterion separately'', \citep{brookhart2013create}. 

\textbf{Generality level} details the descriptions for the criteria, and refers to whether prompts are \textit{general} or \textit{task-specific} (i.e., evaluation prompts that [can$\backslash$cannot] also be used for other tasks) \citep{brookhart2013create}. Autorater evaluation prompts can be viewed through this lens - for example, a holistic autorater evaluation prompt may ask for a single overall judgment whereas an analytic prompt would decompose the evaluation into criterion to be evaluated separately.


% a holistic rubric embedding many interrelated sub-criteria can be more complex than a simple analytic one.


\textbf{Prompt complexity}, as used in this work, refers to the cognitive demands placed on a rater during evaluation. Prior work on task complexity and cognitive demand has identified the number of \textit{paths}, or components, that must be considered simultaneously \citep{8163de9b-fb12-3776-8b6f-4c14075da748}, the degree of element interactivity (i.e., the extent to which components must be processed together rather than independently) \citep{sweller_element_2010} and ambiguity resulting in communication failure \citep{8163de9b-fb12-3776-8b6f-4c14075da748} as characteristics of cognitive load. Thus in our work, we consider the number of criteria evaluated, the degree to which score-level descriptions are interrelated, and the extent to which the rater must resolve ambiguity across sub-criteria. 

\subsection{LLM-as-judges}
% \jessica{talk about the stuff reviewer B said, what have other people done that's not actually the stat analysis}

The autorater (LLM-as-a-judge) paradigm has received substantial recent attention. Several past works demonstrated that LLMs can produce evaluations consistent with human experts. \citep{chiang_can_2023, li-etal-2025-generation}. However, research also shown that autoraters are sensitive to prompt variations, with different instructions leading to substantially different performances \cite{mizrahi-etal-2024-state} and varying quality, necessitating statistical procedures to justify replacing human raters \cite{calderon-etal-2025-alternative}. Several recent works have studied how prompt modifications affect autorater performance. \citep{sclar2024quantifying} demonstrated that prompt formatting choices (e.g., separator characters, whitespace) can significantly affect task accuracy, though these effects weakly correlate across models. \citep{lu-etal-2022-fantastically} showed that example ordering in few-shot prompts substantially impacts performance on classification tasks. Our work extends this literature by using statistical procedures where possible to test how rubric modification choices, such as decomposition level, example selection, and aggregation methods affect agreement with human judgment on tasks where even trained human raters often disagree. This work also empirically examines whether simpler rubrics, often assumed to reduce cognitive load, actually improve human-autorater agreement, in domains where there is no single correct answer, contrasting with the classification and multiple-choice tasks used in prior work.

% Several recent works have studied how prompt modifications affect LLM performance. \citealp{sclar2024quantifying} demonstrated that prompt formatting choices (e.g., separator characters, whitespace) can significantly affect task accuracy, though these effects weakly correlate across models. \citep{lu-etal-2022-fantastically} showed that example ordering in few-shot prompts substantially impacts performance on classification tasks. Both works use paired statistical tests (e.g., chi-squared) to assess whether one model or prompt configuration outperforms another on tasks with known correct answers. \jessica{can we get a bit here on how examples help humans, a citation? i think we mention it later on specifically for aes but doing like a precursor while we're talking about this could be good cuz we mention rubrics that were made for autoraters perform better}

% \alfredo{The following paragraph may be trimmed since parts are restated in earlier paragraph}
% This work differs in a few ways. First, this work examines the interactions between multiple rubric modification components (decomposition level, example selection, and domain characteristics) rather than modifying single rubric elements. It expands on prior studies by examining whether simpler prompts, which are often assumed to reduce cognitive load, improve human-autorater agreement. In addition, the domains considered are ones in which there is no single correct answer, where trained human raters often do not fully agree, contrasting with the classification and multiple choice tasks used in prior work. This work uses comparisons against human judgment rather than a fixed ground truth.

\subsection{Automatic Essay Scoring (AES) and Instruction-Following (IF)}

Recent work has explored using autoraters for AES, employing various personas, including a ``virtual evaluator with expertise in English composition'' \citep{10.1145/3706468.3706507}, a ``helpful pattern-following assistant'' \citep{mansour-etal-2024-large}, and an ``English essay writing test evaluation committee'' or ``English teacher'' \citep{lee-etal-2024-unleashing}. However, these personas may be misaligned with the original human raters; for example, scoring guidelines for a portion of the dataset from \citet{asap-aes} explicitly state that raters should not be teachers. The number of few-shot examples provided to the autorater also varies, \citet{10.1145/3706468.3706507} select the closest three examples with calculated embeddings, while \citet{kundu2024large} chose one essay that scored highly and one essay that scored poorly. 

While most studies utilize the rubrics given by the AES datasets, \citet{lee-etal-2024-unleashing} automatically decomposed the original essay scoring rubric into sub-criteria and performed a modified average aggregation on the sub-criteria, which on average performs better than using a single score from zero-shot prompting. However, \citet{10.1145/3706468.3706507} demonstrated significant improvements by fine-tuning GPT-3.5-turbo and Llama3-8B compared to using GPT-4 with rubrics and few-shot examples. These previous works also may request explanations for scores from autoraters to mimic CoT, aiming for more accurate explanations and ratings. 
% However, studies on the alignment of human and LLM reasoning for open-ended responses, while not for AES, have shown that LLMs exhibit limited correctness and lack  alignment with the top five human responses in a Reddit community \citep{amirizaniani2024can, amirizaniani2024assessing}. 
Studies on analytic rubrics primarily focus on cross-prompt scoring \citep{chen-li-2023-pmaes}. The experimental setup in this work follows the rigor of prior work by 
studying multiple autoraters, using given rubrics, and performing various rubric edits. However, the hypotheses examined extend prior findings on edited rubrics and further investigate rubric components across various essay scoring rubrics.

Research on instruction-following in LLMs has led to two primary approaches for improving and evaluating LLMs, instruction tuning and alignment tuning. Prior work has found that larger models tend to follow instructions more accurately, though this relationship is not strictly linear \citep{ouyang2022training}. \citet{honovich-etal-2023-instruction} proposed allowing LLMs to write instructions based on only seeing examples of a task, although this is still less accurate than human-written instructions. 


% \alfredo{This section seems repetitive now, should cut}
% Unlike these works, which focus on direct evaluation, this study investigates how instruction decomposition influences the alignment between human and model judgments. \citet{sclar2024quantifying} found that prompt formatting weakly correlates across models, reinforcing sensitivity to instruction presentation and the importance of careful instruction design. This work expands on prior studies by examining whether simpler prompts, which are often assumed to reduce cognitive load, enhance alignment between humans and autoraters.

% \jessica{bring back simpler prompts}